\documentclass[11pt,a4paper]{article}

% ---------------------------------------------------------------
% Paquetes
% ---------------------------------------------------------------
\usepackage[utf8]{inputenc}
\usepackage[T1]{fontenc}
\usepackage[spanish,es-tabla,es-nodecimaldot]{babel}
\usepackage{lmodern}
\usepackage[margin=2.5cm]{geometry}
\usepackage{booktabs}
\usepackage{tabularx}
\usepackage{longtable}
\usepackage{array}
\usepackage{amsmath}
\usepackage{xcolor}
\usepackage{listings}
\usepackage{enumitem}
\usepackage{float}
\usepackage[hidelinks]{hyperref}

\definecolor{codebg}{RGB}{246,248,250}
\definecolor{codekw}{RGB}{0,92,197}
\definecolor{codecm}{RGB}{106,115,125}
\definecolor{codestr}{RGB}{3,47,98}
\definecolor{ok}{RGB}{26,127,55}
\definecolor{warn}{RGB}{154,103,0}

\lstset{
  basicstyle=\ttfamily\small,
  backgroundcolor=\color{codebg},
  keywordstyle=\color{codekw}\bfseries,
  commentstyle=\color{codecm}\itshape,
  stringstyle=\color{codestr},
  frame=single,
  rulecolor=\color{black!15},
  breaklines=true,
  columns=fullflexible,
  showstringspaces=false,
  xleftmargin=4pt,
  xrightmargin=4pt,
}
\lstdefinelanguage{shell}{morekeywords={mvn,docker,git,cd,curl},sensitive=true,morecomment=[l]\#}

\newcommand{\ok}{\textcolor{ok}{\textbf{Cumple}}}
\newcommand{\parcial}{\textcolor{warn}{\textbf{Cumple con observación}}}
\setlength{\emergencystretch}{3em} % evita que nombres de clases en \texttt salgan del margen
\newcolumntype{L}{>{\raggedright\arraybackslash}X}

% ---------------------------------------------------------------
% Datos de la entrega (completar)
% ---------------------------------------------------------------
\newcommand{\curso}{Integración Continua}
\newcommand{\entrega}{Entrega 1 -- Pruebas automatizadas}
\newcommand{\integrantes}{%
  Integrante A \\ Integrante B \\ Integrante C \\ Integrante D}
\newcommand{\docente}{Nombre del docente}
\newcommand{\repositorio}{\url{https://github.com/Velasquez076/orders-ci}}

\title{\textbf{\entrega}\\[4pt]\large Cálculo del total de pedidos con descuentos y validación de existencias}
\author{\integrantes\\[6pt]\small \curso\ -- Semestre 7\\ \small Docente: \docente}
\date{8 de noviembre de 2026}

\begin{document}
\maketitle

\begin{abstract}
Este documento presenta la aplicación mínima construida para representar una parte del sistema de
pedidos de la cadena de tiendas de barrio: el cálculo del total de un pedido con descuentos y la
validación de existencias contra una base de datos. Se describe cómo se definió primero el
comportamiento esperado, las 25 pruebas automáticas que lo verifican (ejecutables con un solo
comando), la prueba que reproduce el error de descuentos del caso y las razones por las que la
automatización reduce el costo y el tiempo frente a la revisión manual mensual. También se explica,
con transparencia, una limitación del historial de versiones respecto al orden en que se escribieron
pruebas y código.
\end{abstract}

\setcounter{tocdepth}{1}
\tableofcontents
\newpage

% ===============================================================
\section{Resumen de cumplimiento}
% ===============================================================

\begin{table}[H]
\centering
\begin{tabularx}{\textwidth}{@{}L L l@{}}
\toprule
\textbf{Requisito} & \textbf{Evidencia} & \textbf{Estado}\\
\midrule
Aplicación mínima: total con descuentos y validación de existencias &
Java 21 + Spring Boot 3 (sección~\ref{sec:app}) & \ok\\
Al menos cinco pruebas automáticas & 25 pruebas (sección~\ref{sec:pruebas}) & \ok\\
Ejecutables con un solo comando & \texttt{mvn test} o \texttt{./scripts/evidencias.sh docker} & \ok\\
Prueba que reproduce el error de descuentos del caso &
\texttt{CP04} (sección~\ref{sec:cp04}) & \ok\\
Comportamiento esperado definido primero, con entradas y resultados &
\texttt{docs/ESPECIFICACION.md} (sección~\ref{sec:spec}) & \ok\\
Pruebas escritas antes que el código & Proceso y limitación del historial (sección~\ref{sec:tdd}) & \parcial\\
Por qué la automatización reduce costo y tiempo & Sección~\ref{sec:costos} & \ok\\
\bottomrule
\end{tabularx}
\caption{Requisitos de la entrega y dónde se evidencian.}
\end{table}

Repositorio: \repositorio.

% ===============================================================
\section{Contexto del problema}
% ===============================================================

La empresa es una cadena de tiendas de barrio que gestiona pedidos y despachos. En el caso descrito,
los descuentos llegaron mal a los clientes: el porcentaje se aplicaba \emph{dos veces} sobre el mismo
pedido, de modo que el cliente pagaba menos de lo debido y la empresa perdía margen en cada venta con
descuento. El control existente era una \textbf{revisión manual mensual}, por lo que un error de este
tipo podía permanecer activo hasta un mes antes de ser detectado, y su corrección dependía de que una
persona revisara a mano pedidos, totales y existencias.

El objetivo de esta entrega es mostrar que esas mismas verificaciones pueden expresarse como pruebas
automáticas que se ejecutan en minutos, cada vez que el código cambia.

% ===============================================================
\section{Aplicación}
\label{sec:app}
% ===============================================================

Se eligió \textbf{Java 21} con \textbf{Spring Boot 3.3}, \textbf{Maven} como herramienta de
construcción, \textbf{JUnit 5} y \textbf{AssertJ} para las pruebas, y \textbf{PostgreSQL} como base
de datos (H2 en memoria durante las pruebas).

\begin{table}[H]
\centering
\begin{tabularx}{\textwidth}{@{}l L@{}}
\toprule
\textbf{Componente} & \textbf{Responsabilidad}\\
\midrule
\texttt{OrderCalculator} & Subtotal, descuento (una sola vez sobre el subtotal), total y redondeo a 2 decimales.\\
\texttt{InventoryService} & Consulta y descuento de existencias en la tabla \texttt{inventory}.\\
\texttt{OrderService} & Caso de uso: valida existencias de todos los ítems y calcula el total.\\
\texttt{OrderController} & \texttt{POST /api/orders/total}: responde 200, 400 (datos inválidos) o 409 (sin existencias).\\
\texttt{InfoController} & \texttt{GET /api/info}: indica el ambiente activo (\texttt{dev}, \texttt{test}, \texttt{prod}).\\
\bottomrule
\end{tabularx}
\caption{Componentes principales.}
\end{table}

La aplicación se ejecuta en Docker como aplicación multicontenedor (API + PostgreSQL) mediante
\texttt{docker compose}, y las pruebas pueden ejecutarse en un contenedor independiente sin necesidad
de instalar Java ni Maven.

% ===============================================================
\section{Comportamiento esperado (definido antes de programar)}
\label{sec:spec}
% ===============================================================

Antes de escribir código se definieron las reglas de negocio y una tabla de casos con su entrada y su
resultado esperado. Este documento es el contrato que las pruebas verifican
(\texttt{docs/ESPECIFICACION.md}).

\subsection{Reglas de negocio}

\begin{table}[H]
\centering
\begin{tabularx}{\textwidth}{@{}l L@{}}
\toprule
\textbf{Id} & \textbf{Regla}\\
\midrule
R1 & Subtotal = suma de (cantidad $\times$ precio unitario) de cada ítem.\\
R2 & El descuento es un porcentaje entre 0 y 100 y se aplica \textbf{una sola vez}, sobre el subtotal del pedido.\\
R3 & Total = subtotal $-$ descuento. Nunca es negativo.\\
R4 & Los valores se redondean a 2 decimales (HALF\_UP).\\
R5 & Un pedido debe tener al menos un ítem; cantidad $> 0$; precio $\geq 0$.\\
R6 & Un pedido solo se acepta si hay existencias suficientes de cada SKU (se suman las cantidades del mismo SKU).\\
\bottomrule
\end{tabularx}
\caption{Reglas de negocio.}
\end{table}

\subsection{Casos esperados}

\begin{longtable}{@{}l p{6.6cm} p{5.6cm}@{}}
\toprule
\textbf{Caso} & \textbf{Entrada} & \textbf{Resultado esperado}\\
\midrule
\endfirsthead
\toprule
\textbf{Caso} & \textbf{Entrada} & \textbf{Resultado esperado}\\
\midrule
\endhead
CP01 & $2 \times 10.000 + 3 \times 5.000$ & Subtotal 35.000\\
CP02 & $3 \times 10.000$, descuento 0\,\% & Total 30.000,00\\
CP03 & $1 \times 100.000$, descuento 10\,\% & Total 90.000,00\\
\textbf{CP04} & $2 \times 50.000 + 3 \times 20.000$, descuento 10\,\% & \textbf{Total 144.000,00} (no 129.600)\\
CP05 & $2 \times 10.000$, descuento 100\,\% & Total 0,00\\
CP06 & Descuento 100,01\,\% & Error de validación\\
CP07 & Descuento $-5$\,\% & Error de validación\\
CP08 & $3 \times 3.333{,}33$, descuento 15\,\% & Total 8.499,99\\
CP09 & Pedido sin ítems & Error de validación\\
IN01 & Existencias 10, piden 4 & Hay existencias\\
IN02 & Existencias 10, piden 10 & Hay existencias\\
IN03 & Existencias 10, piden 11 & No hay existencias\\
IN04 & SKU inexistente & No hay existencias\\
IN05 & Descontar 3 de 10 & Quedan 7\\
IN06 & Descontar 11 de 10 & Error; el inventario no cambia\\
PS01 & Caso CP04 con existencias suficientes & Subtotal 160.000, descuento 16.000, total 144.000\\
PS02 & Piden 6, existen 5 & Error de existencias\\
PS03 & Mismo SKU en dos líneas (3 + 3), existen 5 & Error de existencias\\
PS04 & Procesar un pedido & El inventario no cambia\\
WE01--03 & \texttt{POST /api/orders/total} & 200 con total / 409 sin existencias / 400 datos inválidos\\
\bottomrule
\caption{Casos de prueba definidos antes de la implementación.}
\end{longtable}

% ===============================================================
\section{Pruebas automatizadas}
\label{sec:pruebas}
% ===============================================================

\subsection{Inventario de pruebas}

Cada caso de la especificación tiene una prueba con el mismo identificador en su nombre, lo que
permite trazar requisito $\rightarrow$ prueba $\rightarrow$ resultado.

\begin{table}[H]
\centering
\begin{tabularx}{\textwidth}{@{}l L r@{}}
\toprule
\textbf{Clase de prueba} & \textbf{Qué verifica} & \textbf{Pruebas}\\
\midrule
\texttt{OrderCalculatorTest} & CP01--CP09: subtotal, descuento, redondeo y validaciones & 9\\
\texttt{InventoryServiceTest} & IN01--IN06: existencias contra base de datos (H2) & 6\\
\texttt{OrderServiceTest} & PS01--PS04: caso de uso completo de un pedido & 4\\
\texttt{OrderControllerTest} & WE01--WE03: respuestas HTTP 200, 409 y 400 & 3\\
\texttt{OpenApiDocumentationTest} & Contrato OpenAPI y disponibilidad de Swagger UI & 2\\
\texttt{OrdersApplicationTests} & La aplicación arranca correctamente & 1\\
\midrule
\multicolumn{2}{@{}l}{\textbf{Total}} & \textbf{25}\\
\bottomrule
\end{tabularx}
\caption{Pruebas automáticas por clase.}
\end{table}

\subsection{Ejecución con un solo comando}

Con Java 21 y Maven instalados:

\begin{lstlisting}[language=shell]
mvn test
\end{lstlisting}

Sin instalar Java ni Maven, dentro de un contenedor Docker (guarda además la salida con fecha en
\texttt{evidencias/05\_docker\_tests.txt}):

\begin{lstlisting}[language=shell]
./scripts/evidencias.sh docker
\end{lstlisting}

Las pruebas también se ejecutan automáticamente en GitHub Actions en cada \emph{push} y en cada
\emph{pull request} (\texttt{.github/workflows/ci.yml}).

\subsection{Resultado}

Ejecución del 8 de noviembre de 2026 en contenedor (Docker Desktop, contexto \texttt{desktop-linux}):

\begin{lstlisting}
[INFO] Tests run: 25, Failures: 0, Errors: 0, Skipped: 0
[INFO] BUILD SUCCESS
\end{lstlisting}

% ===============================================================
\section{Prueba que reproduce el error de descuentos del caso}
\label{sec:cp04}
% ===============================================================

El error del caso consiste en aplicar el descuento \emph{dos veces} al mismo pedido. Con el pedido
del caso CP04 la diferencia es:

\begin{align*}
\text{Subtotal} &= 2 \times 50.000 + 3 \times 20.000 = 160.000\\
\text{Total correcto (descuento una vez)} &= 160.000 \times (1 - 0{,}10) = 144.000\\
\text{Total con el error (descuento dos veces)} &= 160.000 \times (1 - 0{,}10)^2 = 129.600
\end{align*}

Es decir, el error regala 14.400 adicionales (9\,\% del subtotal) en un solo pedido. La prueba fija
el valor correcto y, explícitamente, rechaza el valor que produce el defecto:

\begin{lstlisting}[language=Java]
@Test
void CP04_discountIsNotAppliedTwiceToTheSameOrder_reproducesCaseBug() {
    List<OrderItem> items = List.of(item("RICE", 2, "50000"), item("OIL", 3, "20000"));
    // subtotal = 160000 ; 10 % = 16000 ; correct total = 144000

    BigDecimal total = calculator.calculateTotal(items, new BigDecimal("10"));

    assertThat(total).isEqualByComparingTo("144000.00");
    // Result produced by the defect (double discount): 129600
    assertThat(total).isNotEqualByComparingTo("129600.00");
}
\end{lstlisting}

Si alguien vuelve a introducir el doble descuento (por ejemplo, aplicándolo por ítem y luego sobre el
total), esta prueba falla en el siguiente \texttt{mvn test} y el pipeline de integración continua
bloquea el cambio antes de que llegue a producción. El mismo escenario se verifica de extremo a
extremo en \texttt{PS01} (caso de uso) y \texttt{WE01} (API REST).

% ===============================================================
\section{Desarrollo guiado por pruebas y limitación del historial}
\label{sec:tdd}
% ===============================================================

\subsection{Proceso seguido}

El trabajo siguió el orden del desarrollo guiado por pruebas (TDD):

\begin{enumerate}[noitemsep]
  \item Se definió el comportamiento esperado: reglas R1--R6 y la tabla de casos con entradas y
        resultados (\texttt{docs/ESPECIFICACION.md}).
  \item Se escribieron las pruebas a partir de esa tabla, por ejemplo
        \texttt{OrderCalculatorTest} e \texttt{InventoryServiceTest}. Al no existir aún la implementación, las pruebas
        fallaban (fase \emph{roja}).
  \item Se escribió el código mínimo que las hace pasar (fase \emph{verde}): la calculadora
        (\texttt{OrderCalculator}), el servicio de existencias, el caso de uso y los controladores.
  \item Se mejoró el código manteniendo todas las pruebas en verde (refactorización).
\end{enumerate}

Esta trazabilidad queda reflejada en el propio código: cada prueba lleva el identificador del caso de
la especificación (CP01, IN03, PS02, \ldots), y la clase \texttt{OrderCalculatorTest} documenta que fue
escrita antes de la implementación.

\subsection{Limitación: el historial de Git no refleja ese orden}

La guía del equipo (\texttt{docs/GUIA\_GIT\_Y\_EVIDENCIAS.md}) planificaba subir primero la
especificación, después las pruebas (\emph{rojo}) y luego la implementación (\emph{verde}) en commits
separados, de forma que el historial demostrara el orden. Sin embargo, el proyecto se publicó en el
repositorio remoto en \textbf{un único commit} (\texttt{359db62}, ``add project to remote
repository''), que contiene especificación, pruebas y código a la vez. Por eso el historial de Git
\textbf{no permite verificar} por sí solo que las pruebas se escribieron antes que el código.

Se decidió \textbf{no reescribir el historial} de la rama principal: ya está publicada y es compartida
por el equipo, y reescribirla exigiría un \emph{force-push} que podría hacer perder trabajo de otros
integrantes. En su lugar, se deja constancia aquí de la limitación. La evidencia disponible del orden
de trabajo es:

\begin{itemize}[noitemsep]
  \item La especificación con entradas y resultados esperados, independiente del código.
  \item La correspondencia uno a uno entre los casos de la especificación y los nombres de las pruebas.
  \item La prueba \texttt{CP04}, construida a partir del error descrito en el caso y no a partir de la
        implementación.
\end{itemize}

Como lección aprendida, en las siguientes entregas cada cambio se trabajará en una rama propia
(\texttt{feature/*}, como ya se hizo con \texttt{feature/docker}) y las pruebas se subirán en un
commit anterior al de la implementación.

% ===============================================================
\section{Por qué la automatización reduce costo y tiempo}
\label{sec:costos}
% ===============================================================

\subsection{Revisión manual mensual frente a pruebas automáticas}

\begin{table}[H]
\centering
\begin{tabularx}{\textwidth}{@{}l L L@{}}
\toprule
\textbf{Aspecto} & \textbf{Revisión manual mensual} & \textbf{Pruebas automáticas}\\
\midrule
Momento de detección & Hasta 30 días después de que el error entra en producción & Minutos después del cambio, antes de llegar a producción\\
Duración de una revisión & Horas o días de una persona revisando pedidos y totales & Segundos: 25 pruebas en una sola ejecución\\
Costo por revisión & Se paga cada mes en horas de personal & Casi nulo tras escribir la prueba una vez\\
Cobertura & Una muestra de pedidos; depende de la atención del revisor & Los mismos casos, completos, en cada ejecución\\
Repetibilidad & Varía según la persona y el día & Idéntica cada vez; resultado objetivo (pasa / falla)\\
Errores que vuelven & Pueden reaparecer sin que nadie lo note & \texttt{CP04} impide que el doble descuento regrese\\
Evidencia & Informes manuales & Historial de ejecuciones en el pipeline de CI\\
\bottomrule
\end{tabularx}
\caption{Comparación entre la revisión manual mensual y las pruebas automáticas.}
\end{table}

\subsection{Costo de un error que se detecta tarde}

El costo de un defecto como el del caso crece con el tiempo que permanece sin detectar. De forma
general:

\[
\text{Pérdida} \approx P \times d \times \bar{s} \times e
\]

donde $P$ es el número de pedidos con descuento por día, $d$ los días hasta que se detecta el error,
$\bar{s}$ el subtotal promedio de esos pedidos y $e$ la fracción perdida por pedido (9\,\% del
subtotal en el caso del doble descuento del 10\,\%).

Con revisión mensual, $d$ puede llegar a 30 días; con pruebas automáticas, el error no llega a
producción y $d \approx 0$. Por ejemplo, \emph{con valores ilustrativos} de 50 pedidos con descuento
al día y un subtotal promedio de 160.000, la pérdida en 30 días sería
$50 \times 30 \times 160.000 \times 0{,}09 = 21.600.000$, frente a cero si la prueba bloquea el cambio.
\textit{(Reemplazar por las cifras reales de la empresa si el caso las proporciona.)}

\subsection{Razones para esta empresa}

\begin{enumerate}
  \item \textbf{El error se detecta antes de que llegue a los clientes.} Cada \emph{push} ejecuta las
        pruebas en GitHub Actions; un cambio que rompe el cálculo del descuento no se integra.
  \item \textbf{Se elimina el costo recurrente de la revisión.} Las horas que hoy se dedican cada mes
        a revisar a mano pedidos y existencias pueden dedicarse a otras tareas; escribir la prueba
        cuesta una sola vez.
  \item \textbf{Se reduce el costo de corregir.} Un error detectado minutos después de escribirlo se
        corrige en el mismo momento, por la persona que lo introdujo y con el contexto fresco. Uno
        detectado un mes después exige rastrear qué cambió, corregir pedidos ya facturados y, a veces,
        dar explicaciones a los clientes.
  \item \textbf{La verificación es completa y objetiva.} Las pruebas comprueban siempre todos los
        casos definidos (descuentos, redondeo, límites, existencias) en lugar de una muestra, y no
        dependen del cansancio o la experiencia del revisor.
  \item \textbf{Permite cambiar el sistema con confianza.} La cadena puede agregar nuevas promociones
        o tiendas: si un cambio afecta un comportamiento existente, una prueba lo señala de inmediato.
  \item \textbf{Las existencias también se protegen.} Las pruebas IN01--IN06 y PS02--PS03 evitan que
        se acepten pedidos sin inventario suficiente, lo que evitaría despachos fallidos y clientes
        insatisfechos.
\end{enumerate}

% ===============================================================
\section{Conclusiones}
% ===============================================================

\begin{itemize}
  \item La aplicación cubre el cálculo del total con descuentos y la validación de existencias contra
        una base de datos, con 25 pruebas automáticas que pasan y se ejecutan con un solo comando.
  \item La prueba \texttt{CP04} reproduce el error de doble descuento del caso y garantiza que no vuelva
        a producción.
  \item El comportamiento se definió primero en una especificación con entradas y resultados
        esperados; el historial de Git no refleja el orden pruebas-antes-que-código, y se documenta
        esa limitación en lugar de ocultarla.
  \item Frente a la revisión manual mensual, la automatización detecta los errores en minutos en vez
        de semanas, elimina un costo recurrente y ofrece una verificación completa y repetible.
\end{itemize}

\end{document}
